% ============================================================================
% Core Feature: Document Management
% ============================================================================
\section{Document Management}

\begin{ugModule}{Overview}

  \textbf{Purpose:}
  Upload, organize, search, and manage PDF documents with rich
  metadata, multi-level approval routing, automatic text
  extraction, and thumbnail generation.

  \textbf{Who Can Access:}
  \ugRole{Admin} \ugRole{Manager} \ugRole{Staff}
\end{ugModule}

% ────────────────────────────────────────────────────────────────
\subsection{Document List Page}

The Document List page (\ugMenu{Documents > Document List}) is
the primary view for browsing all documents you have access to.

\subsubsection*{List Columns}

The data grid displays the following columns:

\begin{tabularx}{\textwidth}{@{} l l X @{}}
  \toprule
  \textbf{Column} & \textbf{Sortable}
    & \textbf{Description} \\
  \midrule
  No
    & No
    & Sequential row number \\
  \ugField{Document Title}
    & Yes
    & Document title with an access indicator icon
      showing the access level: Public
      (\,\faIcon{globe}\,), Protected
      (\,\faIcon{shield-alt}\,), or Private
      (\,\faIcon{lock}\,) \\
  \ugField{Year}
    & Yes
    & Fiscal/calendar year of the document
      (sorted by \texttt{year.code}) \\
  \ugField{Unit}
    & Yes
    & Requesting organizational unit (sorted by
      \texttt{requestingUnit.name}) \\
  \ugField{Status}
    & Yes
    & Current workflow status displayed as a
      multi-step progress track showing all
      approval levels. If rejected, shows the
      rejection reason and date \\
  \ugField{Issued Date}
    & Yes
    & Date the document was officially issued
      (formatted as day/month/year) \\
  \ugField{Submitted Date}
    & Yes
    & Date and time the document was submitted
      (includes time) \\
  Action
    & No
    & Context menu with actions: View Detail, Edit,
      Delete \\
  \bottomrule
\end{tabularx}

\begin{ugTip}
  Click any sortable column header to toggle between ascending
  and descending order. A small arrow indicator shows the
  current sort direction.
\end{ugTip}

% ────────────────────────────────────────────────────────────────
\subsection{Creating a New Document}

Creating a document is a multi-step process using a
\textbf{tabbed wizard} with three tabs:

\begin{center}
  \fbox{\textbf{Tab 1:} Basic Information}\quad
  $\longrightarrow$\quad
  \fbox{\textbf{Tab 2:} Upload File}\quad
  $\longrightarrow$\quad
  \fbox{\textbf{Tab 3:} Access Control}
\end{center}

\vspace{6pt}

Two fields are always visible \textbf{above the tabs}:

\begin{tabularx}{\textwidth}{@{} l l l X @{}}
  \toprule
  \textbf{Field} & \textbf{Type} & \textbf{Required}
    & \textbf{Description} \\
  \midrule
  \ugField{Document Title}
    & Text input
    & \textbf{Yes}
    & The official title of the document. Max 1024
      characters. This is the primary identifier shown
      in lists, search results, and detail pages. \\
  \ugField{Year}
    & Dropdown
    & \textbf{Yes}
    & Select the fiscal or calendar year that the
      document belongs to. Populated from the Year
      master data. \\
  \bottomrule
\end{tabularx}

% ────────────────────────────────────────────────────────────────
\subsection{Tab 1: Basic Information}

This tab captures the document's metadata and classification.
Fields are grouped into a card layout.

\subsubsection*{Core Fields (Always Visible)}

\begin{tabularx}{\textwidth}{@{} l l l X @{}}
  \toprule
  \textbf{Field} & \textbf{Type} & \textbf{Required}
    & \textbf{Description} \\
  \midrule
  \ugField{Requesting Unit}
    & Dropdown
    & \textbf{Yes}
    & The organizational unit submitting or owning the
      document (e.g., IT, HR, Finance, Legal). Loaded
      from the Unit master data. Used for filtering and
      reporting. \\
  \ugField{Approval Purpose}
    & Dropdown
    & \textbf{Yes}
    & Defines the approval routing for this document.
      Each purpose maps to a preconfigured chain of
      up to 6 approval levels, determining who must
      sign off on the document. The dropdown is
      filtered by the selected Unit. \\
  \ugField{Description}
    & Text area
    & \textbf{Yes}
    & A summary of the document's content and purpose.
      Max 4096 characters. Displayed on the document
      detail page and included in search indexing. \\
  \ugField{Submitted Date}
    & Date picker
    & \textbf{Yes}
    & The date the document is submitted to the system.
      Auto-filled with today's date and locked
      (non-editable). \\
  \bottomrule
\end{tabularx}

\subsubsection*{Conditional Date Fields}

The following date fields appear \textbf{only if enabled} by
the selected Document Type's configuration. Each Document Type
can independently toggle which dates are relevant:

\begin{tabularx}{\textwidth}{@{} l l l X @{}}
  \toprule
  \textbf{Field} & \textbf{Type} & \textbf{Condition}
    & \textbf{Description} \\
  \midrule
  \ugField{Issued Date}
    & Date picker
    & \texttt{hasIssuedDate}
    & The date the document was officially issued or
      signed by the authority. \\
  \ugField{Expired Date}
    & Date picker
    & \texttt{hasExpiredDate}
    & The date after which the document is no longer
      valid. After this date, access is automatically
      denied. \\
  \ugField{Effective Date}
    & Date picker
    & \texttt{hasEffectiveDate}
    & The date when the document takes effect (becomes
      enforceable). \\
  \ugField{Publish Date}
    & Date picker
    & \texttt{hasPublishDate}
    & The date the document was published or made
      available to its audience. \\
  \bottomrule
\end{tabularx}

\begin{ugNote}
  The Document Type controls which date fields appear on the
  form. For example, an ``SK'' (Surat Keputusan) may require
  Issued Date and Effective Date, while a ``Proposal'' may
  only need Submitted Date. This is configured by the Admin
  in the Document Type master data.
\end{ugNote}

\subsubsection*{JDIH Fields (Conditional)}

If the selected Document Type has \texttt{isJDIH = true}
(legal document repository classification), three additional
fields appear:

\begin{tabularx}{\textwidth}{@{} l l l X @{}}
  \toprule
  \textbf{Field} & \textbf{Type} & \textbf{Required}
    & \textbf{Description} \\
  \midrule
  \ugField{Agency}
    & Dropdown
    & No
    & The government or corporate agency that issued
      the document. Loaded from the Agency master data. \\
  \ugField{JDIH Category}
    & Dropdown
    & No
    & Legal document classification category according
      to JDIH standards (e.g., Peraturan, Keputusan,
      Instruksi). \\
  \ugField{Document Source}
    & Dropdown
    & No
    & The origin or provenance of the document (e.g.,
      Internal, External, Government). \\
  \bottomrule
\end{tabularx}

% ────────────────────────────────────────────────────────────────
\subsection{Tab 2: Upload File}

This tab provides a file upload area using a
\textbf{drag-and-drop dropzone}:

\begin{tabularx}{\textwidth}{@{} l l l X @{}}
  \toprule
  \textbf{Field} & \textbf{Type} & \textbf{Required}
    & \textbf{Description} \\
  \midrule
  \ugField{Document File}
    & File dropzone
    & \textbf{Yes}
    & The PDF file to upload. You can drag and drop a
      file onto the dropzone area, or click the zone
      to open a file browser dialog. \\
  \bottomrule
\end{tabularx}

\subsubsection*{File Validation Rules}

\begin{tabularx}{\textwidth}{@{} l X @{}}
  \toprule
  \textbf{Rule} & \textbf{Details} \\
  \midrule
  Accepted formats
    & \texttt{application/pdf} (\texttt{.pdf}) only.
      Other file types are rejected immediately. \\
  MIME type check
    & The file's Content-Type header must be
      \texttt{application/pdf}. \\
  Extension check
    & The file must have a \texttt{.pdf} extension. \\
  Magic bytes check
    & The first bytes of the file must contain
      \texttt{\%PDF} (the PDF file signature). This
      prevents disguised files from being uploaded. \\
  Maximum file size
    & Configurable per Document Type (default: 50~MB
      global). The Document Type master data allows
      setting a custom \ugField{Max File Size (MB)} per
      type. \\
  \bottomrule
\end{tabularx}

\begin{ugWarning}
  Only single-file uploads are supported. Each document
  record corresponds to exactly one PDF file. To manage
  multiple related files, create separate documents and
  use tags to group them.
\end{ugWarning}

% ────────────────────────────────────────────────────────────────
\subsection{Tab 3: Access Control}

This tab configures who can view and download the document.
Access is controlled using an \textbf{ACL (Access Control List)}
toggle group:

\subsubsection*{Access Levels}

\begin{tabularx}{\textwidth}{@{} l l X @{}}
  \toprule
  \textbf{Level} & \textbf{Icon}
    & \textbf{Description} \\
  \midrule
  \ugField{Public}
    & \faIcon{globe}
    & Anyone with a system account can view and
      download this document. No additional access
      configuration needed. \\
  \ugField{Protected}
    & \faIcon{shield-alt}
    & All authenticated users within the organization
      can access, but external sharing is restricted. \\
  \ugField{Private}
    & \faIcon{lock}
    & Only specifically granted users, roles, or units
      can access. You must configure at least one
      access grant. \\
  \bottomrule
\end{tabularx}

\subsubsection*{Private Access Configuration}

When \ugField{Private} is selected, a sub-panel appears with
tabs for three access grant types:

\begin{tabularx}{\textwidth}{@{} l X @{}}
  \toprule
  \textbf{Access Type} & \textbf{Description} \\
  \midrule
  \textbf{By Unit}
    & Grant access to all users within selected
      organizational units. Select one or more units
      from the data grid. \\
  \textbf{By Role}
    & Grant access by unit division level (position).
      Select from a data grid of unit division roles. \\
  \textbf{By User}
    & Grant access to specific individual users. Select
      from a searchable user data grid. \\
  \bottomrule
\end{tabularx}

\begin{ugImportant}
  When \ugField{Private} access is selected, you must grant
  access to \textbf{at least one} unit, role, or user. The
  form will not submit without this. This prevents accidentally
  creating documents that no one can access.
\end{ugImportant}

\begin{ugTip}
  The default access level is pre-set by the Document Type
  configuration (\texttt{defaultAclId}). For example, all
  ``SOP'' documents may default to \ugField{Protected}, while
  ``SK'' documents default to \ugField{Public}. You can
  override this on a per-document basis.
\end{ugTip}

% ────────────────────────────────────────────────────────────────
\subsection{Step-by-Step: Creating a Document}

\begin{ugSteps}
  \item Navigate to \ugMenu{Documents > Document List}.
  \item Click \ugButton{+ Create Document} in the toolbar.
  \item Enter the \ugField{Document Title} and select the
        \ugField{Year} (these are always visible above the
        tabs).
  \item \textbf{Tab 1 --- Basic Information:}
    \begin{itemize}[leftmargin=14pt]
      \item Select the \ugField{Requesting Unit}.
      \item Select the \ugField{Approval Purpose} (this defines
            the approval chain).
      \item Enter a \ugField{Description}.
      \item Fill in any visible date fields (Issued, Expired,
            Effective, Publish) as needed.
      \item If JDIH fields appear, fill in Agency, Category,
            and Source as applicable.
    \end{itemize}
  \item \textbf{Tab 2 --- Upload File:}
    \begin{itemize}[leftmargin=14pt]
      \item Drag and drop your PDF file onto the dropzone, or
            click to browse.
      \item Wait for the file to be accepted (the dropzone
            shows the file name and size).
    \end{itemize}
  \item \textbf{Tab 3 --- Access Control:}
    \begin{itemize}[leftmargin=14pt]
      \item Select Public, Protected, or Private.
      \item If Private: grant access by Unit, Role, or User.
    \end{itemize}
  \item Click \ugButton{Save} to submit.
\end{ugSteps}

% ────────────────────────────────────────────────────────────────
\subsection{What Happens After Upload}

When you click \ugButton{Save}, the system processes the
document through the following pipeline:

\begin{enumerate}[leftmargin=18pt]
  \item The PDF file is \textbf{uploaded} to a temporary
        storage location first, then associated with the
        document record.
  \item \textbf{All three forms are validated} simultaneously
        (Info, Upload, Access). If any form has errors, you
        are notified and the submission stops.
  \item The file is \textbf{validated} --- MIME type, extension,
        magic bytes (\texttt{\%PDF}), and file size.
  \item A \textbf{SHA-256 checksum} is computed for integrity
        verification.
  \item The file is stored in \textbf{MinIO} object storage at
        path \texttt{documents/\{id\}/\{version\}/\{filename\}}.
  \item A \textbf{restricted (rasterized) copy} is generated
        and stored at the restricted file path for secure
        downloads.
  \item \textbf{Text is extracted} from all pages for full-text
        search indexing.
  \item A \textbf{thumbnail} is generated from the first page
        and stored in MinIO at
        \texttt{thumbnails/\{id\}.bmp}.
  \item The document is created with version~1 in
        \ugStatus{ugMutedText}{Draft} status.
  \item The document's \textbf{approval chain} is initialized
        based on the selected Approval Purpose.
  \item An \texttt{AI Ingest} event is published via Redis
        Streams to asynchronously generate vector embeddings
        for the RAG search engine.
\end{enumerate}

% ────────────────────────────────────────────────────────────────
\subsection{Document Detail Page}

After creating or opening a document, the detail page displays
the following sections:

\subsubsection*{General Information Card}

Displays key metadata in a read-only format:

\begin{tabularx}{\textwidth}{@{} l X @{}}
  \toprule
  \textbf{Field} & \textbf{Description} \\
  \midrule
  Document Title
    & The document's title \\
  Agency
    & Issuing agency (shown only for JDIH documents) \\
  JDIH Category
    & Legal classification (shown only for JDIH
      documents) \\
  Document Source
    & Origin of the document (shown only for JDIH
      documents) \\
  Requesting Unit
    & Organizational unit that owns the document \\
  Year
    & Fiscal/calendar year \\
  Description
    & Full text description \\
  Access Level
    & Current ACL level (Public, Protected, Private)
      with corresponding icon \\
  \bottomrule
\end{tabularx}

\subsubsection*{Document Preview Card}

Displays an embedded PDF preview using an iframe:

\begin{itemize}[leftmargin=18pt]
  \item The PDF is rendered in a 3:4 aspect ratio container
  \item Two action buttons: \ugButton{View PDF} (opens in
        new tab) and \ugButton{Download}
  \item If no file is uploaded, a placeholder message is shown
\end{itemize}

\subsubsection*{Statistics Cards}

Four stat cards displayed in a horizontal row:

\begin{tabularx}{\textwidth}{@{} l l X @{}}
  \toprule
  \textbf{Stat} & \textbf{Icon}
    & \textbf{Description} \\
  \midrule
  Total Views
    & \faIcon{eye}
    & Number of times the document was previewed \\
  Total Downloads
    & \faIcon{download}
    & Number of times the document was downloaded \\
  Unique Users
    & \faIcon{users}
    & Count of distinct users who accessed it \\
  Last Access
    & \faIcon{clock}
    & Timestamp of the most recent access event \\
  \bottomrule
\end{tabularx}

\subsubsection*{Charts \& Activity}

\begin{itemize}[leftmargin=18pt]
  \item \textbf{Download Trend Chart} --- Line chart showing
        download volume over time
  \item \textbf{View Trend Chart} --- Line chart showing
        preview volume over time
  \item \textbf{Activity Timeline} --- Chronological list of
        all document events (created, updated, downloaded,
        status changed, access granted)
\end{itemize}

% ────────────────────────────────────────────────────────────────
\subsection{Updating a Document}

To update a document's metadata without changing the file:

\begin{ugSteps}
  \item Navigate to the document detail page.
  \item Click \ugButton{Edit} from the action menu.
  \item The same 3-tab form appears, pre-filled with existing
        data.
  \item Modify any fields across all three tabs.
  \item On the Upload tab, you can either keep the existing
        file or upload a new one (creating a new version).
  \item Click \ugButton{Save} to apply changes.
\end{ugSteps}

\begin{ugNote}
  Metadata-only updates (title, description, unit, dates)
  do \textbf{not} create a new file version. Only uploading
  a new PDF file increments the version number.
\end{ugNote}

% ────────────────────────────────────────────────────────────────
\subsection{Document Type Configuration}

The behavior of the document form is driven by the
\textbf{Document Type} master data. Administrators configure
each type with the following options:

\begin{tabularx}{\textwidth}{@{} l l X @{}}
  \toprule
  \textbf{Setting} & \textbf{Type}
    & \textbf{Effect} \\
  \midrule
  \texttt{isJDIH}
    & Boolean
    & Shows Agency, JDIH Category, and Document
      Source fields on the form \\
  \texttt{isCert}
    & Boolean
    & Marks as certification document (special
      handling) \\
  \texttt{isISO}
    & Boolean
    & Marks as ISO-related document \\
  \texttt{hasIssuedDate}
    & Boolean
    & Shows Issued Date field \\
  \texttt{hasExpiredDate}
    & Boolean
    & Shows Expired Date field \\
  \texttt{hasEffectiveDate}
    & Boolean
    & Shows Effective Date field \\
  \texttt{hasPublishDate}
    & Boolean
    & Shows Publish Date field \\
  \texttt{requiresApproval}
    & Boolean
    & Determines if the document must go through the
      approval workflow \\
  \texttt{defaultAclId}
    & Reference
    & Pre-selects the default access level \\
  \texttt{maxFileSizeMb}
    & Integer
    & Overrides the global max file size for this
      type \\
  \texttt{retentionDays}
    & Integer
    & Number of days before automatic retention
      expiry \\
  \texttt{notifyBeforeExpireDays}
    & Integer
    & Number of days before expiry to trigger
      notification alerts \\
  \texttt{allowedExtensions}
    & Text
    & Comma-separated list of allowed file
      extensions \\
  Watermark Access
    & Boolean
    & Include access-level watermark on downloads \\
  Watermark Confidential
    & Boolean
    & Include confidentiality watermark on downloads \\
  Watermark Company
    & Boolean
    & Include company branding watermark on
      downloads \\
  \bottomrule
\end{tabularx}

\begin{ugTip}
  Document Types allow the system to adapt its behavior
  per document category. For example, a ``Surat Keputusan''
  type might require Issued Date and Effective Date, enable
  all watermarks, and default to Public access --- while a
  ``Proposal'' type might only need Submitted Date and
  default to Private access.
\end{ugTip}

% ────────────────────────────────────────────────────────────────
\subsection{Business Rules}

\begin{itemize}[leftmargin=18pt]
  \item Only \texttt{application/pdf} files are accepted.
        The system validates MIME type, file extension, and
        magic bytes (\texttt{\%PDF}).
  \item Maximum file size is configurable per Document Type
        (default: 50~MB global limit).
  \item Documents are created in
        \ugStatus{ugMutedText}{Draft} status with version~1.
  \item The Document Code is auto-generated by the system.
  \item The Owner User is automatically set to the
        logged-in user.
  \item Documents with an \ugField{Expired Date} are
        automatically denied access after that date.
  \item Text extraction and AI ingestion run asynchronously
        and do not block the upload process.
  \item Each document version stores its own SHA-256 checksum,
        original filename, file size, and uploader identity.
  \item Soft-deleted documents (\texttt{IsActive = false})
        are hidden from the list but retained in the database.
\end{itemize}

\newpage
